\chapter{Worked Example on the Toy Curve}
\label{ch:worked}

% ============================================================
%  Full hand computation of kG for k = 1..9 on y^2 = x^3 + 7 mod 17
% ============================================================

\section{Setup}

\[
  p = 17, \quad a = 0, \quad b = 7, \quad G = (1, 5), \quad k = 6 .
\]
Check: $5^2 = 25 \equiv 8$ and $1^3 + 7 = 8$, so $G \in E$. \checkmark

\section{Repeated addition: $G, 2G, \ldots, 9G$}

\begin{framedex}[label={ex:2G}]{$2G$ (doubling)}
  \[
    \lambda = \frac{3 x_1^2 + a}{2 y_1} = \frac{3}{10} = 3 \cdot 10^{-1} = 3 \cdot 12 = 36 \equiv 2
  \]
  \[
    x_3 = 2^2 - 1 - 1 = 2, \qquad y_3 = 2(1 - 2) - 5 = -7 \equiv 10 .
  \]
  So $2G = (2, 10)$. Check: $10^2 = 100 \equiv 15$, $2^3 + 7 = 15$. \checkmark
\end{framedex}

\begin{framedex}{$3G = 2G + G = (2,10) + (1,5)$}
  \[
    \lambda = \frac{5 - 10}{1 - 2} = \frac{-5}{-1} = 5, \qquad
    x_3 = 25 - 2 - 1 = 22 \equiv 5, \qquad
    y_3 = 5(2 - 5) - 10 = -25 \equiv 9 .
  \]
  So $3G = (5, 9)$.
\end{framedex}

\begin{framedex}{$4G = 3G + G = (5,9) + (1,5)$}
  \[
    \lambda = \frac{-4}{-4} = 1, \qquad
    x_3 = 1 - 5 - 1 = -5 \equiv 12, \qquad
    y_3 = 1(5 - 12) - 9 = -16 \equiv 1 .
  \]
  So $4G = (12, 1)$.
\end{framedex}

\begin{framedex}{$5G = 4G + G = (12,1) + (1,5)$}
  \[
    \lambda = \frac{4}{-11} = \frac{4}{6} = 4 \cdot 6^{-1} = 4 \cdot 3 = 12,
  \]
  \[
    x_3 = 144 - 12 - 1 = 131 \equiv 12, \qquad
    y_3 = 12(12 - 12) - 1 \equiv 16 .
  \]
  So $5G = (12, 16) = -4G$.
\end{framedex}

\begin{framedex}{$6G = 5G + G = (12,16) + (1,5)$}
  \[
    \lambda = \frac{-11}{-11} = 1, \qquad
    x_3 = 1 - 12 - 1 = -12 \equiv 5, \qquad
    y_3 = 1(12 - 5) - 16 = -9 \equiv 8 .
  \]
  So $6G = (5, 8)$: this is the \textbf{toy public key}.
\end{framedex}

Continuing in the same way:

\begin{center}
  \begin{tabular}{c c c c c c}
    \toprule
    $k$ & operation & numerator / denominator & $\lambda$ & $kG$ & note \\
    \midrule
    1 & -- & -- & -- & $(1, 5)$ & $G$ \\
    2 & $G + G$ & $3 / 10$ & 2 & $(2, 10)$ & doubling \\
    3 & $2G + G$ & $-5 / -1$ & 5 & $(5, 9)$ & \\
    4 & $3G + G$ & $-4 / -4$ & 1 & $(12, 1)$ & \\
    5 & $4G + G$ & $4 / -11$ & 12 & $(12, 16)$ & $= -4G$ \\
    6 & $5G + G$ & $-11 / -11$ & 1 & $\mathbf{(5, 8)}$ & $= -3G$, \textbf{public key} \\
    7 & $6G + G$ & $-3 / -4$ & 5 & $(2, 7)$ & $= -2G$ \\
    8 & $7G + G$ & $-2 / -1$ & 2 & $(1, 12)$ & $= -G$ \\
    9 & $8G + G$ & $-7 / 0$ & -- & $\mathcal{O}$ & vertical line \\
    \bottomrule
  \end{tabular}
\end{center}

\begin{framedobs}{What the table tells us}
  \begin{itemize}
    \item At $k = 9$ the denominator $x_2 - x_1$ is $0$: $8G$ and $G$ have the same $x$ and opposite $y$, so $9G = \mathcal{O}$. The \textbf{order of $G$ is 9}, and from now on the sequence repeats: $10G = G$, $11G = 2G$, \ldots
    \item The sequence is symmetric: $(9 - k) G = -kG$, i.e. same $x$, $y$ replaced by $17 - y$. For instance $6G = (5,8)$ and $3G = (5,9)$.
    \item $G$ visits only 8 of the 17 affine points. The other 9 points are never reached from $G$ (they are outside the subgroup $\abk{G}$).
    \item The code must handle the case ``denominator $= 0$'' explicitly (return $\mathcal{O}$), otherwise the modular inverse of 0 is undefined and the program crashes or loops.
  \end{itemize}
\end{framedobs}

\section{Double-and-add: $6G$ with 3 bits}

$6 = (110)_2$. Starting from $R = \mathcal{O}$:
\begin{enumerate}
  \item bit 1: $R = 2\mathcal{O} + G = G = (1,5)$;
  \item bit 1: $R = 2G = (2,10)$, then $R = 2G + G = (5,9)$;
  \item bit 0: $R = 2 \cdot (5,9)$:
  \[
    \lambda = \frac{3 \cdot 5^2}{2 \cdot 9} = \frac{75}{18} \equiv \frac{7}{1} = 7, \qquad
    x_3 = 49 - 10 = 39 \equiv 5, \qquad
    y_3 = 7(5 - 5) - 9 \equiv 8 .
  \]
\end{enumerate}
Result: $6G = (5, 8)$, the same as with repeated addition. \checkmark

\begin{framedex}{Expected output of the program}
\begin{lstlisting}[numbers=none]
Curve: y^2 = x^3 + 7 (mod 17)
Base point G = (1, 5)   order of G = 9
Private key k = 6

1G = (1, 5)
2G = (2, 10)
3G = (5, 9)
4G = (12, 1)
5G = (12, 16)
6G = (5, 8)

TOY PUBLIC KEY: Q = 6G = (5, 8)\end{lstlisting}
\end{framedex}

\section{Program structure (JavaScript)}

\begin{lstlisting}[language=Java]
const P = 17n, A = 0n, B = 7n;          // toy curve, BigInt
const G = { x: 1n, y: 5n };
const O = null;                          // point at infinity

const mod = (a) => ((a % P) + P) % P;
const inv = (a) => modPow(mod(a), P - 2n); // Fermat

function add(p1, p2) {
  if (p1 === O) return p2;
  if (p2 === O) return p1;
  if (p1.x === p2.x && mod(p1.y + p2.y) === 0n) return O;
  const lambda = (p1.x === p2.x && p1.y === p2.y)
    ? mod(3n * p1.x * p1.x + A) * inv(2n * p1.y)   // doubling
    : mod(p2.y - p1.y) * inv(p2.x - p1.x);         // addition
  const l  = mod(lambda);
  const x3 = mod(l * l - p1.x - p2.x);
  const y3 = mod(l * (p1.x - x3) - p1.y);
  return { x: x3, y: y3 };
}

function scalarMul(k, pt) {                // double-and-add
  let R = O;
  for (const bit of k.toString(2)) {
    R = add(R, R);
    if (bit === "1") R = add(R, pt);
  }
  return R;
}\end{lstlisting}
(The language \texttt{Java} is used only for syntax highlighting; \texttt{modPow} is the usual square-and-multiply.) Using \texttt{BigInt} from the start means that the same code works unchanged with the real \texttt{secp256k1} parameters.
